\documentclass[conference]{IEEEtran}
\IEEEoverridecommandlockouts
% The preceding line is only needed to identify funding in the first footnote. If that is unneeded, please comment it out.
\usepackage{cite}
\usepackage{amsmath,amssymb,amsfonts}
\usepackage{algorithmic}
\usepackage{graphicx}
\usepackage{textcomp}
\usepackage{xcolor}
\def\BibTeX{{\rm B\kern-.05em{\sc i\kern-.025em b}\kern-.08em
    T\kern-.1667em\lower.7ex\hbox{E}\kern-.125emX}}

% ---------------------------------------------------------------
% Values marked TODO are not confirmed. Fill from your own outputs.
% ---------------------------------------------------------------
\newcommand{\NTrain}{12{,}705}      % train responses after val holdout
\newcommand{\NVal}{2{,}237}        % val responses
\newcommand{\NValSrc}{377}     % val source IDs (15% of train sources)
\newcommand{\GridSeed}{7}    % dev-grid seed: my notes conflict (7 vs 42), check run_all.sbatch
\newcommand{\GPU}{A100}         % GPU model on hpc3
\newcommand{\LenFeature}{characters}  % unit of the length feature in baselines.py (tokens, words, chars)

% ---------------------------------------------------------------
% Every value I could not confirm prints as ?? in bold.
% Fill each \tbd from your script output, never from memory.
% ---------------------------------------------------------------
\newcommand{\tbd}{\textbf{??}}

\begin{document}

\title{Grounding Improves Response-Level Hallucination Detection Unevenly Across Task Families: A ModernBERT Input Ablation on RAGTruth}

\author{\IEEEauthorblockN{Oleksandr Pometun}
\IEEEauthorblockA{\textit{Institute of Cognitive Science} \\
\textit{University of Osnabrueck}\\
Osnabrueck, Germany \\
opometun@uni-osnabrueck.de}
}

\maketitle

\begin{abstract}
Hallucination detectors on RAGTruth use the retrieved source, but how much
their performance depends on it has not been measured. We fine-tune
ModernBERT-base for response-level detection with and without the source.
Backbone, data, training and selection are identical, with three seeds per
condition. Adding the source raises positive-class $F_1$ from 0.671 to 0.711,
a mean gain of 0.040, concentrated in summarization and question answering and
absent in data-to-text. Logistic regression on task family, response length and
generator identity reaches 0.700, above the blind detector, so RAGTruth's
response-level labels are predictable to a considerable degree from metadata
alone. Grounding
helps, but this test set limits what blind performance can show.
\end{abstract}

\begin{IEEEkeywords}
hallucination detection, retrieval-augmented generation, ModernBERT, RAGTruth, input ablation
\end{IEEEkeywords}

\section{Introduction}

Retrieval-augmented generation gives a language model retrieved documents, 
allowing it to base its answers on evidence from them\cite{lewis2020rag}. 
However, the generator may produce statements not supported by the retrieved text, 
and these hallucinations can be hard to detect because they sound convincing \cite{ji2023survey}.
A RAG system in production often requires a detector that monitors and
flags unsupported responses. This detector should be cheap enough to
be used on every generated response.

RAGTruth \cite{niu2024ragtruth} contains roughly 18k responses from six generators
across three task categories (question answering, summarization and data-to-text),
with hallucinated text passages labeled at the span level. On this dataset,
encoder detectors compete with much larger language models. LettuceDetect
\cite{kovacs2025lettucedetect} fine-tunes ModernBERT
\cite{warner2024modernbert}, a BERT-style encoder \cite{devlin2019bert}
supporting an 8{,}192-token context, for token classification across source,
question and answer. Its large version achieves an example-level F1 of 79.22\%,
while its base version, using the same ModernBERT-base backbone as ours,
achieves 76.07\%.

Although all these detectors use the source, previous work has not measured how
much their performance relies on it instead of patterns in the response related
to task, length and generator style. Natural language inference demonstrates the
same risk: models that receive only the hypothesis, without the premise, perform
above chance by using annotation cues
\cite{gururangan2018annotation,poliak2018hypothesis}. For hallucination detection,
a source-blind detector serves as the equivalent baseline.

We adapt the ModernBERT-on-RAGTruth configuration from
\cite{kovacs2025lettucedetect} for response-level classification and evaluate it
using a controlled change to the input. Condition~A receives the task
instruction and response, whereas Condition~B additionally receives the
retrieved source. We investigate (i)~the performance
gain from adding the source, (ii)~how this gain varies across task families, and
(iii)~how Condition~A performs against baselines based on task family, response
length and generator identity. Because published results use different setups,
we treat them only as reference points, with Condition~A as our experimental
baseline.

\section{Methods}

\subsection{Data and Labels}
We follow the official RAGTruth split and remove the same 173 quality-flagged
responses from both conditions (144 incorrect refusals and 29 truncated
responses). The resulting dataset contains 17{,}617 responses, including
2{,}675 in the test set. RAGTruth has no validation split, so we use 15\% of the training source IDs
(\NValSrc{} sources) for validation while preserving the task-family
distribution. This produces \NTrain{} training responses and \NVal{} validation
responses. A source can have up to six responses, one from each generator, so
we split by source instead of response to keep all responses from a source in
the same partition. The code checks that source IDs do not overlap. We use the same split for all
conditions, configurations and seeds. A response is hallucinated if it has at
least one annotated span, $y=\mathbf{1}[\,|S(r)|\ge 1\,]$, where $S(r)$ is the
set of annotated hallucination spans of response $r$.
Unsupported content counts as hallucinated even when factually true (RAGTruth's
\texttt{implicit\_true} spans are kept). Positive prevalence is about 0.45 in
training and validation and 0.35 in test.

\subsection{Inputs}
Let $t$ denote the fixed instruction for each task family, including the
question for question answering; let $r$ denote the response and $s$ the source
text (\texttt{source\_info}). Each leaf in a structured source is 
written as one \texttt{key: value} line.
The input segments are then joined in the following order:
\begin{align}
x_A &= \texttt{[CLS]}\; t\; r\; \texttt{[SEP]}, \label{eq:xa}\\
x_B &= \texttt{[CLS]}\; t\; r\; \texttt{[SEP]}\; s\; \texttt{[SEP]}. \label{eq:xb}
\end{align}
We build $x_A$ without the source, and an automated check confirms that changing
the source alters every $x_B$ and no $x_A$. We limit sequences to
$T_{\max}=4096$ tokens and allow only $s$ to be truncated. Because the longest
input contains 2{,}632 tokens, no input is truncated in either condition.

\subsection{Model}
We use ModernBERT-base (22 layers and hidden size $d=768$) as the encoder
$f_\theta$ \cite{warner2024modernbert}. Its context window can hold every source
in full, unlike a 512-token encoder. Given an input of $T$ tokens and an
attention mask $m\in\{0,1\}^T$, let $M=\sum_{i=1}^{T} m_i$ be the number of real
(non-padding) tokens. Then
\begin{align}
H &= f_\theta(x) \in \mathbb{R}^{T\times d}, \label{eq:enc}\\
\bar{h} &= \frac{1}{M}\sum_{i=1}^{T} m_i H_i \in \mathbb{R}^{d}, \label{eq:pool}\\
z &= W\bar{h} + b, \qquad W\in\mathbb{R}^{2\times d},\; b\in\mathbb{R}^{2}. \label{eq:head}
\end{align}
Our head is a single newly initialized linear layer with no dropout,
not the stock ModernBERT classification head. The prediction is positive (hallucinated) when
$z_1 \ge z_0$, with ties assigned to the positive class.

\subsection{Training and Selection}
We fine-tune all parameters using a custom PyTorch loop that minimizes
cross-entropy loss over a batch of $N$ examples,
\begin{equation}
\mathcal{L} = -\frac{1}{N}\sum_{n=1}^{N} \log p_{n,y_n},
\qquad p_n = \mathrm{softmax}(z_n),
\label{eq:loss}
\end{equation}
using AdamW \cite{loshchilov2019adamw} with weight decay 0.01 (none on bias and
normalization parameters), gradient-norm clipping at 1.0 
and gradient accumulation to achieve an effective batch size of 16. 
The learning rate increases linearly for the first 10\% of update steps and
then decreases linearly to zero. Training runs for at most
six epochs. Validation macro-F1 is computed after every epoch, training stops
after two epochs without improvement, and the best checkpoint is restored.

The maximum learning rate $\eta$ is the only hyperparameter we tune. Using seed
\GridSeed{} and validation data only, we compare $10^{-5}$, as used by
\cite{kovacs2025lettucedetect}, with $3\times10^{-5}$ under both conditions. We
select $\eta=3\times10^{-5}$ because it gives the higher mean validation
macro-F1 across A and B (0.758 versus 0.746). We then fix this configuration for
six final runs: three seeds for each condition, all run on one \GPU{} GPU of the
university HPC cluster. The
seeds determine the classifier-head initialization and data order. 
We use them to measure training variation, not as statistical pairs.
The backbone, examples, data split, optimizer, schedule and selection rule are
the same for both conditions. Only the input length and, because of early
stopping, the selected epoch necessarily differ. We use the test set only to
evaluate the six final checkpoints.

\subsection{Evaluation and Diagnostics}
The primary metric is positive-class F1, $F_1^{+}=2PR/(P+R)$, with precision
$P$ and recall $R$ of the hallucinated class computed over the whole test set.
We also report macro-F1, the mean of $F_1^{+}$ and the clean-class $F_1^{-}$,
and the per-seed gain $\Delta = F_{1,B}^{+}-F_{1,A}^{+}$. All scores use the
rule $z_1 \ge z_0$ above; threshold-free analysis is left to future work.

We use three reference baselines to place Condition~A in context. 
If a fraction $p$ of responses are hallucinated, always predicting
hallucinated gives $F_1^{+}=2p/(1+p)$; a real model can still score below this.
 The second baseline predicts the majority
class for each task family, using the training labels to determine that class.
The third trains logistic regression on the training and validation data using
one-hot task family and response length (\LenFeature{}) as features.
An additional version also includes generator identity, a dataset
attribute that may be unavailable to a deployed detector. These baselines show
how much of Condition~A's score these features could explain, but they do not
prove what the model has learned.

To see whether Condition~B uses the source, we test its three trained models
again on a modified test set. Within each task family, every response receives
the source of another response, never its own. We use the same fixed
reassignment for all three models. Only the source changes, the instruction and
response remain the same. For responses that were originally clean, we report
the predicted-positive rate before and after the swap and the fraction of
predictions that change.

\section{Results and Discussion}

All scores are on the 2{,}675 test responses (943 hallucinated, prevalence
0.353). Table~\ref{tab:systems} compares the systems and
Table~\ref{tab:family} breaks the comparison down by task family.

\begin{table}[t]
\caption{Test results. P, R and $F_1^{+}$ refer to the hallucinated class.
A and B are means over three seeds. Mean $F_1^{+}$ is averaged over runs, not
recomputed from mean P and R. $q_+$ is the predicted-positive rate among the 1{,}732 originally clean
responses. $F_1^{+}$ seed ranges are [0.667, 0.675] for A and [0.708, 0.714]
for B.}
\label{tab:systems}
\centering
\footnotesize
\setlength{\tabcolsep}{3pt}
\begin{tabular}{lccccc}
\hline
System & P & R & $F_1^{+}$ & Macro-$F_1$ & $q_+$ \\
\hline
Always positive            & 0.353 & 1.000 & 0.521 & 0.261 & 1.00 \\
Majority per family        & 0.643 & 0.614 & 0.628 & 0.716 & 0.19 \\
Metadata (family, length)$^{\mathrm{a}}$ & 0.640 & 0.616 & 0.628 & 0.715 & 0.19 \\
A: source-blind            & 0.671 & 0.670 & 0.671 & 0.746 & 0.18 \\
B: grounded                & 0.710 & 0.715 & 0.711 & 0.776 & 0.16 \\
B, substituted source$^{\mathrm{b}}$ & -- & -- & -- & -- & 0.98 \\
\hline
\multicolumn{6}{l}{$^{\mathrm{a}}$Adding generator identity as a feature gives $F_1^{+}=0.700$.}\\
\multicolumn{6}{l}{$^{\mathrm{b}}$Source-substitution diagnostic; only $q_+$ is reported.}
\end{tabular}
\end{table}
% NOTE: if you prefer, replace the P/R/F1/Macro cells of the last row with
% dashes. The FPR cell is the readout that matters there.

\begin{table}[t]
\caption{Task-family breakdown (means over three seeds). $n$ and prevalence
are shared by A and B. Flip is the share of B's predictions on clean responses
that change under source substitution.}
\label{tab:family}
\centering
\setlength{\tabcolsep}{4pt}
\begin{tabular}{lcccccc}
\hline
Family & $n$ & Prev. & A $F_1^{+}$ & B $F_1^{+}$ & $\Delta$ & Flip \\
\hline
Summary   & 900 & 0.227 & 0.198 & 0.414 & $+0.217$ & 0.89 \\
QA        & 875 & 0.183 & 0.456 & 0.567 & $+0.111$ & 0.85 \\
Data2txt  & 900 & 0.643 & 0.850 & 0.840 & $-0.009$ & 0.62 \\
\hline
\end{tabular}
\end{table}

\textit{How large is the gain?} Supplying the source raises $F_1^{+}$ from
0.671 to 0.711 (Table~\ref{tab:systems}). The per-seed gains are $+0.035$,
$+0.041$ and $+0.044$ (mean $+0.040$), and the seed ranges of A and B do not
overlap on $F_1^{+}$ or on macro-$F_1$. With three runs per condition this
shows a consistent direction, not a powered estimate of the effect size.
% TODO (one sentence, needs your P and R): say whether B gains through
% precision, recall or both, e.g. "The gain comes mainly through ...".

\textit{Where is the gain concentrated?} The overall result hides large
differences across task families (Table~\ref{tab:family}). For summarization,
Condition~A scores $F_1^{+}=0.198$, below the always-positive reference for that
family (0.370), and one A run collapses to 0.028. Condition~B reaches 0.414,
with per-seed gains of 0.11, 0.11 and 0.43. The gain is $+0.111$ for question answering, whereas
the source provides no improvement for data-to-text ($-0.009$). We consider two
possible explanations for this pattern, but the current experiments cannot
distinguish between them. First, a summary must be compared with its article.
From the response alone, a fluent summary of another article appears equally
valid. In contrast, data-to-text responses contain surface cues that the blind
detector may exploit. Second, hallucination prevalence varies across task
families. The high prevalence in data-to-text increases positive-class $F_1$
for any detector biased toward positive predictions. Condition~A's macro-$F_1$
of 0.772 for this family shows that it still distinguishes between the classes.
However, summarization scores remain low even for Condition~B (0.414). One
possible reason is our pooling method. When source tokens dominate the input,
mean pooling may weaken the signal from a short unsupported span. A token-level
objective, as used by \cite{kovacs2025lettucedetect}, may be more suitable for
such cases.

\textit{How much of Condition~A can metadata baselines explain?} Logistic
regression using task family and response length achieves
$F_1^{+}=0.628$. This equals the majority-per-family rule: the regression
predicts the positive class for all data-to-text responses and almost no others,
so response length adds nothing beyond task family.
Condition~A scores only slightly higher at 0.671, while both
outperform the always-positive baseline of 0.521. Adding generator identity
raises the metadata baseline to 0.700, above Condition~A. This weakens what
detector performance on this test set can tell us about unseen generators. We therefore make no transfer claims for
either condition.

\textit{Does Condition~B use the source?} The paired comparison shows the effect of adding the source, but not how
Condition~B uses it. The source-substitution test provides further evidence.
When we pair each test response with another source from the same task family, the
predicted-positive rate across the 1{,}732 clean responses increases from 0.16
to 0.98. In addition, 82\% of their predictions change, with consistent results
across all three Condition~B models. Condition~B is therefore highly sensitive
to the source it receives. However, the substituted pairs were not relabeled,
so this diagnostic cannot measure correctness. After substitution the predicted-positive rate on
originally clean responses is at least 0.97 in every family, so the lower flip
rate on data-to-text (0.62) only reflects its higher starting rate. The
near-complete change points to a broad mismatch signal more than claim-level
checking.

\textit{Optimization and overfitting.} By validation macro-$F_1$, early stopping
selected epochs 3, 3 and 1 for Condition~A: it reached its best validation score
early, and further epochs brought no improvement. Condition~B selected epochs 4,
6 and 5, all three hitting the six-epoch cap and one peaking in the final epoch,
so it may not have fully converged and the reported gain holds for this training
budget.

\textit{Limitations and trade-offs.} Both conditions share one learning rate,
selected on the mean validation score of A and B. This is a fairness choice,
and B's longer inputs might prefer a different value. Three seeds per condition
give direction, not power. The cues available to A are tied to this dataset's
generators and tasks. Finally, grounding has a cost: B's inputs are much longer
than A's (mean 797 against 183 tokens),
so the gain of 0.040 is bought with a corresponding increase in inference
compute.

\section{Conclusion}

Supplying the source raises $F_1^{+}$ from 0.671 to 0.711 (mean 0.040), with
gains in summarization and question answering and none in data-to-text.
Condition~A modestly exceeds a family-and-length baseline and falls below its
generator-aware version, so we make no transfer claims. Source substitution
shows sensitivity to source-response mismatch, not correct verification, and
grounding also requires much longer inputs (mean 797 against 183 tokens).
Threshold-free evaluation, a source-level bootstrap and a token-level objective
are the next steps.

\section*{Acknowledgment}
Funded by the Deutsche Forschungsgemeinschaft (DFG, German Research Foundation) -- 456666331.

\bibliographystyle{IEEEtran}
\bibliography{refs}

\end{document}
